\section{Noise Analysis}
\label{sec:results:noise}
This section presents the noise growth of the implemented functions, reported throughout as the max-norm $\|\epsilon\|_\infty = \max_i\{\epsilon_i\}$. The absolute magnitude of the error is of little interest on its own; what matters is how many bits of the modulus it occupies. Every plot in this section is therefore semi-logarithmic with a base-2 $y$ axis, so that a vertical distance reads directly as a number of bits.

\subsection{Core Library}
For the external product, we apply $x_i = x_{i-1} \boxdot A_i$ where $x_0 = \Call{Encrypt}{1}$ using fresh encryptions $A_i = \Call{RGSW}{1}$ each iteration.

\subsubsection{External Product}
Due to the asymmetry discussed in~\autoref{rem:extprod-asymmetry} making the noise additive instead of multiplicative, we expect to see linear noise growth with repeated applications of the external product as long as the \gls{rgsw} remains fresh.

% \definecolor{gaussblue}{RGB}{31,119,180}
% \definecolor{hpsgreen}{RGB}{44,160,44}
% \definecolor{canonorange}{RGB}{255,165,0}
% \definecolor{maxred}{RGB}{255,36,0}
% \definecolor{hpspurple}{RGB}{148,103,189}

% ---------- Compute fits in linear space ----------
% \pgfplotstableread[col sep=comma]{Data/Noise/ExternalProd/bv_small.csv}\bvsmall
% \pgfplotstablecreatecol[linear regression={x=n, y=noise}]{reg}{\bvsmall}
% \edef\smallSlope{\pgfplotstableregressiona}
% \edef\smallIntercept{\pgfplotstableregressionb}

\pgfplotstableread[col sep=comma]{Data/Noise/ExternalProd/bv.csv}\bvlarge
\pgfplotstablecreatecol[linear regression={x=n, y=noise}]{reg}{\bvlarge}
\edef\largeSlope{\pgfplotstableregressiona}
\edef\largeIntercept{\pgfplotstableregressionb}

\begin{figure}[H]
    \centering
    \caption{Accumulated noise from chained external products}
    \label{fig:bv:plot-1}
    \begin{tikzpicture}
    \begin{semilogyaxis}[
        noiseaxis,
        log basis y=2,
        xmin=0, xmax=1000,
        xlabel={Chain length},
        ylabel={$\log\|\epsilon\|_\infty$},
        legend pos=north west,
        set layers,
        width=\textwidth,
        height=0.55\textwidth,
    ]
        % ---------- Scatter plots ----------
        % \addplot[only marks, mark=*, mark size=1pt, color=gaussblue]
        %     table[col sep=comma, x=n, y=noise] {Data/Noise/ExternalProd/bv_small.csv};
        \addplot[only marks, mark=*, mark size=1pt, color=hpsgreen]
            table[col sep=comma, x=n, y=noise] {Data/Noise/ExternalProd/bv.csv};
            
        % ---------- Regression lines ----------
        % \addplot[thick, solid, dotted, on layer=axis foreground, color=canonorange, domain=1:1000, samples=300, forget plot]
        %     {\smallSlope*x + \smallIntercept};
        \addplot[thick, solid, dotted, on layer=axis foreground, color=hpspurple, domain=1:1000, samples=300, forget plot]
            {\largeSlope*x + \largeIntercept};
            
        % \legend{Small Parameters, Large Parameters}
        \legend{BV-RNS}
    \end{semilogyaxis}
    \end{tikzpicture}
\end{figure}


% \newpage

% \begin{tabular}{lc} 
%     \toprule
%     \textbf{n} & \textbf{noise} \\ % & \textbf{Role} \\ % Manual Table Header
%     \midrule
    
%     % Reads CSV, maps column headers to macros, and loops through rows
%     \csvreader[head to column names]{Data/Noise/ExternalProd/bv.csv}{}
%     {\n & \noise \\} 
    
%     % \bottomrule
% \end{tabular}


The measured growth is linear in the length of the chain, and the error gains only a few bits across the full chain at every value of $\ell$. Increasing $\ell$ shifts the whole series downwards without changing its slope.

\subsubsection{Internal product}
For the internal product, we apply $X_i = X_{i-1} \boxtimes A_i$ where $X_0 = \Call{RGSW}{1}$ using fresh encryptions $A_i = \Call{RGSW}{1}$ each iteration.

% \item Maybe: show max chain length if RHS is accumulated noise
\definecolor{plotblue}{HTML}{00ADB5}
% \definecolor{hpsgreen}{RGB}{44,160,44}
% \definecolor{canonorange}{RGB}{255,165,0}
% \definecolor{maxred}{RGB}{255,36,0}
% \definecolor{hpspurple}{RGB}{148,103,189}

% ---------- Compute fits in linear space ----------
% \pgfplotstableread[col sep=comma]{Data/Noise/InternalProd/bv_small.csv}\bvsmall
% \pgfplotstablecreatecol[linear regression={x=n, y=noise}]{reg}{\bvsmall}
% \edef\smallSlope{\pgfplotstableregressiona}
% \edef\smallIntercept{\pgfplotstableregressionb}

\pgfplotstableread[col sep=comma]{Data/Noise/InternalProd/bv_2.csv}\bvlarge
\pgfplotstablecreatecol[linear regression={x=n, y=noise}]{reg}{\bvlarge}
\edef\intSlope{\pgfplotstableregressiona}
\edef\intIntercept{\pgfplotstableregressionb}

\begin{figure}[H]
    \centering
    \caption{Accumulated noise from chained internal products.}
    \label{fig:bv:plot-2}
    \begin{tikzpicture}
    \begin{semilogyaxis}[
        noiseaxis,
        log basis y=2,
        xmin=0, xmax=1000,
        xlabel={Chain length},
        ylabel={Error $\|\epsilon\|_\infty$},
        legend pos=south east,
        set layers,
        width=\textwidth,
        height=0.6\textwidth,
    ]
        % ---------- Scatter plots ----------
        % \addplot[only marks, mark=*, mark size=1pt, color=gaussblue]
        %     table[col sep=comma, x=n, y=noise] {Data/Noise/InternalProd/bv_small.csv};
        % \addplot[only marks, mark=*, mark size=1pt, color=gaussblue]
        %     table[col sep=comma, x=n, y=noise] {Data/Noise/InternalProd/bv_1.csv};
        \addplot[only marks, mark=*, mark size=1pt, color=hpsgreen]
            table[col sep=comma, x=n, y=noise] {Data/Noise/InternalProd/bv_2.csv};
        \addplot[only marks, mark=*, mark size=1pt, color=canonorange]
            table[col sep=comma, x=n, y=noise] {Data/Noise/InternalProd/bv_3.csv};
        \addplot[only marks, mark=*, mark size=1pt, color=gaussblue]
            table[col sep=comma, x=n, y=noise] {Data/Noise/InternalProd/bv_4.csv};
            
        % ---------- Regression lines ----------
        % \addplot[thick, solid, dotted, on layer=axis foreground, color=canonorange, domain=1:1000, samples=300, forget plot]
        %     {\smallSlope*x + \smallIntercept};
        % \addplot[thick, solid, dotted, on layer=axis foreground, color=hpspurple, domain=1:1000, samples=300, forget plot]
        %     {\largeSlope*x + \largeIntercept};
        
        % ---------- Decryption failure ---------
        % \addplot[color=maxred, domain=1:1000]{2^120};
            
        \legend{$\ell = 2$ (baseline), $\ell = 3$, $\ell = 4$}
    \end{semilogyaxis}
    \end{tikzpicture}
\end{figure}

Again we see exactly what we expect; the growth is linear in the length of the chain and again shifts downwards with $\ell$. Both the starting error and the amount added per product are substantially larger than for the external product.

\subsection{sPAR}
The protocol was run end-to-end under the multi-party context using OpenFHE's \texttt{NOISE\_FLOODING\_MULTIPARTY} mode. With the settings from \autoref{tab:noise:multiparty}, we were only able to support $\eta = 4$ users. Increasing $\ell$ allows more users, but also degrades the performance. \autoref{tab:noise:multiparty} lists the largest error found in the message matrix $L$ once every client has written. % Decryption fails once the error reaches $\lfloor Q/2 \rfloor$, which is $359$ bits at the parameters of \autoref{tab:bench:params}.

\begin{table}[H]
    \caption{Largest number of error bits in $L$ after 1 full round.}
    \label{tab:noise:multiparty}
    \centering
    \begin{tabular}{cccc}
        \toprule
        $\mathbf{\eta}$ & \textbf{Error bits \\ $(\ell = 2)$} & \textbf{Error bits \\ $(\ell = 3)$} & \textbf{Error bits \\ $(\ell = 4)$} \\
        \midrule
        $2$  & $245$ & $221$ &  \\
        $4$  & $359$ & $288$ &  \\
        $8$  &  ---  & $356$ &  \\
        % $16$ &  ---  &  ---  & \\
        \bottomrule
    \end{tabular}
\end{table}

Even with $\ell = 3$, we were only able to support $\eta = 8$ users. Initially we suspected this was due to the large $t$. Since the $h \gets I_{i,k}\boxtimes\mathsf{notHasWritten}\boxtimes z_{i,d}$ assignment has only a single fresh ciphertext, $z_{i,d}$, a ``dirty'' ciphertext necessarily must appear on the right hand side in one of the internal products, breaking the noise asymmetry from \autoref{rem:extprod-asymmetry}. However, running the same test with $t=2^8$ only increased the headroom by a few bits.

% \definecolor{gaussblue}{RGB}{31,119,180}
% \definecolor{hpsgreen}{RGB}{44,160,44}
% \definecolor{canonorange}{RGB}{255,165,0}
% \definecolor{maxred}{RGB}{255,36,0}
% \definecolor{hpspurple}{RGB}{148,103,189}

% ---------- Compute fits in linear space ----------
% \pgfplotstableread[col sep=comma]{Data/Noise/write-N2.csv}\bvsmall
% \pgfplotstablecreatecol[linear regression={x=n, y=noise}]{reg}{\bvsmall}
% \edef\smallSlope{\pgfplotstableregressiona}
% \edef\smallIntercept{\pgfplotstableregressionb}

\pgfplotstableread[col sep=comma]{Data/Noise/write-N32.csv}\bvlarge
\pgfplotstablecreatecol[linear regression={x=n, y=noise}]{reg}{\bvlarge}
\edef\largeSlope{\pgfplotstableregressiona}
\edef\largeIntercept{\pgfplotstableregressionb}

\begin{figure}[H]
    \centering
    \caption{Noise in $L$ after full round.}
    \label{fig:spar:plot-1}
    \begin{tikzpicture}
    \begin{semilogyaxis}[
        noiseaxis,
        log basis y=2,
        xmin=1.1, xmax=33,
        xlabel={$\eta$},
        ylabel={Error $\|\epsilon\|_\infty$},
        legend pos=south east,
        set layers,
        width=\textwidth,
        height=0.55\textwidth,
    ]
        % ---------- Scatter plots ----------
        % \addplot[only marks, mark=*, mark size=4pt, color=gaussblue]
        %     table[col sep=comma, x expr=\thisrow{n}+1, y=noise] {Data/Noise/write-N2.csv};
        \addplot[only marks, mark=*, mark size=1pt, color=hpsgreen]
            table[col sep=comma, x expr=\thisrow{n}+1, y=noise] {Data/Noise/write-N32.csv};
            
        % ---------- Regression lines ----------
        % \addplot[thick, solid, dotted, on layer=axis foreground, color=canonorange, domain=1:16, samples=16, forget plot]
        %     {\smallSlope*x + \smallIntercept};
        % \addplot[thick, solid, dotted, on layer=axis foreground, color=hpspurple, domain=1:16, samples=16, forget plot]
        %     {\largeSlope*x + \largeIntercept};
            
        \legend{$\ell = 2$} % \legend{$\eta = 2$, $\eta = 16$}
    \end{semilogyaxis}
    \end{tikzpicture}
\end{figure}


% \newpage

% \begin{tabular}{lc} 
%     \toprule
%     \textbf{n} & \textbf{noise} \\ % & \textbf{Role} \\ % Manual Table Header
%     \midrule
    
%     % Reads CSV, maps column headers to macros, and loops through rows
%     \csvreader[head to column names]{Data/Noise/ExternalProd/bv.csv}{}
%     {\n & \noise \\} 
    
%     % \bottomrule
% \end{tabular}


%% TODO: explain difference between the two multi-party modes -> Maybe in preliminaries
%% TODO: justify why the security would need to be evaluated further
Testing without multi-party FHE, where every user share a single secret key, reveals the noise growth from the protocol itself is much smaller; to the point where a practically unbounded number of users (on the order of ${10^{41}}$) is supported. This confirms that the large noise in the multi-party setting comes from the noise flooding rather than the protocol itself. The flooding is intrinsic to distributed decryption, where it statistically hides each party's key share; switching to \texttt{FIXED\_NOISE\_MULTIPARTY}, whose flooding standard deviation is fixed and smaller, should support more users but would weaken that hiding guarantee and require a separate security analysis.